% 割性质的交换论证：T 中 u 到 v 的路径 u-p2-p1-q1-q2-v 经跨割边 f=p1q1；加入轻边 e=uv，删去 f。
\begin{tikzpicture}[>=Stealth,
  v/.style={circle,draw=FigureStroke,fill=FigureBlueFill,line width=.8pt,minimum size=5.5mm,inner sep=0pt},
  t/.style={draw=FigureBlue,line width=1.6pt},
  lab/.style={font=\figurefont\fontsize{7.5}{9}\selectfont},
  note/.style={font=\FigureNoteFont\fontsize{7.5}{9}\selectfont}]
\newcommand{\verts}{%
  \node[v] (p2) at (-.9,.15) {}; \node[v] (p1) at (0,.95) {}; \node[v] (u) at (0,-.75) {$u$};
  \node[v] (q1) at (2.2,.95) {}; \node[v] (q2) at (3.1,.15) {}; \node[v] (v) at (2.2,-.75) {$v$};
  \draw[dashed,FigureMuted,line width=.7pt] (1.1,1.5)--(1.1,-1.3);
  \node[lab] at (-.6,1.35) {$S$}; \node[lab] at (2.9,1.35) {$V\setminus S$};}
\begin{scope}
  \verts
  \draw[t] (u)--(p2)--(p1); \draw[t] (q1)--(q2)--(v);
  \draw[t] (p1)--node[lab,above,text=FigureInk]{$f$}(q1);
  \draw[draw=FigureGreen,line width=1.6pt] (u)--node[lab,below,text=FigureGreen]{$e$}(v);
  \node[lab,align=center] at (1.1,-1.75) {(a) $T+e$含一个环};
\end{scope}
\begin{scope}[xshift=5.2cm]
  \verts
  \draw[t] (u)--(p2)--(p1); \draw[t] (q1)--(q2)--(v);
  \draw[draw=FigureRed,dashed,line width=1.2pt] (p1)--node[lab,above,text=FigureRed]{$f$}(q1);
  \draw[draw=FigureGreen,line width=1.6pt] (u)--node[lab,below,text=FigureGreen]{$e$}(v);
  \node[lab,align=center] at (1.1,-1.75) {(b) $T'=T+e-f$，$w_e\leq w_f$};
\end{scope}
\end{tikzpicture}
